\documentclass[11pt]{article}
\usepackage[margin=1in]{geometry}
\usepackage{amsmath,amssymb,amsthm}
\usepackage{booktabs}
\usepackage{hyperref}
\usepackage{algorithm,algpseudocode}
\usepackage{enumitem}

\newtheorem{theorem}{Theorem}
\newtheorem{proposition}{Proposition}
\newtheorem{corollary}{Corollary}
\newtheorem{definition}{Definition}
\newtheorem{remark}{Remark}
\newtheorem{example}{Example}
\newtheorem{conjecture}{Conjecture}
\newtheorem{assumption}{Assumption}

\title{Visible-Mode Lyapunov Geometry and High-Throughput \\
Spectral Screening of 6F5 Conservative Matrix Fields}
\author{D. Vesterlund}
\date{July 2026}

\begin{document}
\maketitle

\begin{abstract}
We introduce $J^\star$-Lyapunov geometry, a visible-mode framework for arithmetic approximations generated by matrix cocycles. Let $P_n = M_n \cdots M_1$ and $r_n = \varphi(P_n v_0)/\psi(P_n v_0) = p_n/q_n$. We define the denominator mode $a$ as the first Lyapunov mode visible to $\psi$, and the first visible error mode $J^\star$ as the leading mode that changes the projective ratio. Under explicit domination, visibility, and non-cancellation hypotheses, we prove $\delta_{\mathrm{double}}(n) = -\lambda_{J^\star}(n)/\lambda_a(n) + o(1)$. The standard proxy $-\lambda_2/\lambda_1$ is recovered when $a=1$, $J^\star=2$.

We calibrate on Ap\'ery recurrences for $\zeta(2)$ and $\zeta(3)$, where Lyapunov ratios reproduce exact deltas to $\sim 10^{-4}$. We then apply the framework to 6F5 conservative matrix fields. A GPU implementation screens $\sim 10^9$ trajectories per GPU-hour on AMD MI250X, a $\sim 10^4$ raw throughput improvement over exact RNS matrix multiplication. The proxy serves only as a screening statistic; all reported hits are independently verified by exact integer arithmetic. Large-scale experiments identify positive-delta regions, effective-dimension strata, and candidate spectral boundaries underlying the empirical Delta Ladder.

\medskip
\noindent\textbf{Key words:} Lyapunov exponents, irrationality measure, conservative matrix fields, GPU computing, PSLQ, integer relation detection.

\medskip
\noindent\textbf{Data and code availability:} All code, benchmark outputs, and reproducibility scripts are available at \url{https://github.com/VesterlundCoder/deltaproxy} (MIT license).
\end{abstract}

\tableofcontents

\section{Introduction}

The irrationality delta $\delta$ measures the convergence quality of rational approximants $p_n/q_n \to L$. Positive $\delta_\infty > 0$ supports Ap\'ery-style irrationality arguments \cite{apery1979}. Conservative Matrix Fields (CMFs) \cite{raayoni2021,elimelech2024} provide a systematic framework for discovering such approximants, but the search space is astronomical: the 6F5 companion family has $\sim 1.5 \times 10^{15}$ trajectories.

The bottleneck is arithmetic: computing $\delta$ exactly requires multi-hundred-digit integer matrix products whose bit-length grows with depth. A LUMI GPU campaign screening $2.5 \times 10^9$ candidates with RNS arithmetic consumed $\sim$20{,}000 GPU-hours \cite{vesterlund2026}.

We show that the entire cost can be sidestepped. The finite-time Lyapunov spectrum---computed in fixed-width floating-point with QR renormalization---predicts $\delta$ at a cost of $O(d^3 N)$ flops versus $O(d^3 N^2 \log N)$ bit-cost for the exact path. The resulting proxy screens $10^9$ trajectories in under one GPU-hour.

The paper is organized around a single thesis: \emph{the arithmetic convergence quality of a matrix-cocycle-generated approximation is determined not generally by the second Lyapunov exponent, but by the first Lyapunov mode visible to the projective observable.} This theory simultaneously provides the fast screening method.

\subsection{Contributions}

\begin{enumerate}
\item \textbf{Theory:} Definition of arithmetic approximation cocycles, denominator mode $a$, visibility determinant, and first visible error mode $J^\star$.
\item \textbf{Conditional theorem:} $\delta_{\mathrm{double}} = -\lambda_{J^\star}/\lambda_a + o(1)$ under explicit hypotheses, with a proof that handles $n$-dependent Oseledets modes and subexponential visibility coefficients.
\item \textbf{Arithmetic gauge analysis:} Explicit treatment of gauge dependence and normalization conventions for the delta quantity.
\item \textbf{Mechanism:} Exterior-product interpretation of the approximation error.
\item \textbf{Classical calibration:} Validation on Ap\'ery $\zeta(2)$, $\zeta(3)$, and negative-delta PCFs.
\item \textbf{High-dimensional diagnostics:} Full-spectrum proxy $\hat\delta_j = -\lambda_j/\lambda_a$ with gap, degeneracy, and visibility diagnostics.
\item \textbf{High-throughput discovery:} GPU implementation screening $\sim 10^9$ trajectories/GPU-hour, $\sim 10^4$ raw speedup, with matched precision benchmarks (float32 and float64).
\item \textbf{Experimental CMF geometry:} Maps of positive regions, $z$-dependence, effective dimension, Delta Ladder.
\item \textbf{Reproducible verification architecture:} Proxy for screening only; all hits confirmed by exact integer arithmetic and CAS certificates. Complete code at \cite{deltaproxy_repo}.
\end{enumerate}

\section{Arithmetic approximation cocycles}

\begin{definition}[Matrix cocycle]
Let $M_1, M_2, \ldots$ be $d \times d$ matrices and $P_n = M_n \cdots M_1$. Given linear functionals $\varphi, \psi: \mathbb{R}^d \to \mathbb{R}$, the projective observable is
\[
r_n = \frac{\varphi(P_n v_0)}{\psi(P_n v_0)} = \frac{p_n}{q_n}.
\]
\end{definition}

This covers linear recurrences, polynomial continued fractions, Ap\'ery systems, and CMF trajectories.

\subsection{Arithmetic normalization and gauge dependence}
\label{sec:gauge}

The raw delta quantity
\[
\delta_n = -1 - \frac{\log|r_n - r_{2n}|}{\log|q_n|}
\]
is not a purely projective invariant. If we replace $M_n \mapsto c_n M_n$ for nonzero scalars $c_n$, the ratio $r_n$ is unchanged, but the denominator $q_n$ and all Lyapunov exponents are rescaled: $\lambda_i \mapsto \lambda_i + \frac{1}{n}\sum_k \log|c_k|$. The delta quantity therefore depends on the arithmetic gauge.

We adopt the \emph{primitive integer gauge}: all $M_n$ have integer entries in lowest form (no common factor), and $p_n, q_n$ are the resulting integer numerator and denominator. This is the natural gauge for irrationality proofs, since $|q_n|$ is the actual integer denominator whose growth controls the approximation quality.

Other gauges arise naturally:
\begin{itemize}
\item \emph{RNS gauge:} residues modulo a product of primes; $\log|q_n|$ is recovered from the CRT reconstruction.
\item \emph{Denominator-cleared gauge:} $M_n$ scaled so that $\psi(P_n v_0)$ is an integer at each step.
\item \emph{Subexponential equivalence:} two gauges are equivalent if $\sum_k \log|c_k| = o(n\lambda_a)$, i.e.\ the rescaling is subexponential relative to the denominator growth rate.
\end{itemize}

Throughout this paper, all delta values are computed in the primitive integer gauge unless otherwise stated. When comparing Ap\'ery systems, PCFs, and CMFs, we verify that the gauge is consistent by checking that $\log|q_n|$ grows at the rate $n\lambda_a$ predicted by the Lyapunov spectrum.

\subsection{Finite-time Lyapunov exponents: QR vs.\ SVD}
\label{sec:qr_svd}

We distinguish two finite-time spectral quantities:
\[
\lambda_i^{\mathrm{SVD}}(n) = \frac{1}{n}\log\sigma_i(P_n), \qquad
\widehat\lambda_i^{\mathrm{QR}}(n) = \frac{1}{n}\sum_{k=1}^n \log|R_{k,ii}|,
\]
where $\sigma_i(P_n)$ are the singular values of the product matrix and $R_{k,ii}$ are the diagonal entries from per-step QR decomposition ($P_n = Q_n R_n Q_{n-1}^{-1} \cdots$).

The QR-based estimator $\widehat\lambda_i^{\mathrm{QR}}$ converges to the asymptotic Oseledets Lyapunov exponent $\lambda_i$ as $n \to \infty$ under the multiplicative ergodic theorem \cite{oseledets1968}. However, $\widehat\lambda_i^{\mathrm{QR}}(n) \neq \lambda_i^{\mathrm{SVD}}(n)$ in general at finite $n$: the accumulated QR diagonals approximate the singular values only asymptotically.

\begin{remark}
All proxy values reported in this paper are $\widehat\lambda_i^{\mathrm{QR}}(N)$ at depth $N=120$. The calibration in Section~\ref{sec:calibration} validates that the discrepancy $|\widehat\lambda_i^{\mathrm{QR}}(N) - \lambda_i^{\mathrm{SVD}}(N)|$ is negligible at this depth for the systems tested ($< 10^{-3}$ relative error on Ap\'ery recurrences).
\end{remark}

\begin{definition}[Double-depth delta]
\[
\delta_n = -1 - \frac{\log|r_n - r_{2n}|}{\log|q_n|},
\]
using the exact identity $r_n - r_{2n} = (p_n q_{2n} - p_{2n} q_n)/(q_n q_{2n})$.
\end{definition}

\section{Denominator modes and visibility}

\begin{definition}[Denominator mode]
Assume an asymptotic mode expansion $v_n = \sum_i A_i(n) u_i(n)$ where $u_i(n)$ are the Oseledets covariant Lyapunov vectors at time $n$, with $\log|A_i(n)| = n\lambda_i + o(n\lambda_a)$. The denominator mode is
\[
a = \min\{i : \psi(u_i(n)) \neq 0 \text{ for all sufficiently large } n\}.
\]
Then $\log|q_n| = n\lambda_a + o(n\lambda_a)$.
\end{definition}

\begin{definition}[Visibility determinant]
For $j \neq a$, the visibility determinant is
\[
V_{a,j}(\varphi, \psi; n) = \varphi(u_j(n))\psi(u_a(n)) - \varphi(u_a(n))\psi(u_j(n)).
\]
Mode $j$ is \emph{visible} iff $|V_{a,j}(n)|$ is not exponentially small relative to $n\lambda_a$, i.e., $\log|V_{a,j}(n)| = o(n\lambda_a)$.
\end{definition}

\begin{definition}[First visible error mode]
$J^\star$ is the visible mode with the largest exponent $\lambda_j - \lambda_a$:
\[
J^\star = \min\{j \geq 2 : V_{a,j} \not\equiv 0\} \quad \text{(generic case $a=1$)}.
\]
\end{definition}

\section{Visible-mode delta theorem}
\label{sec:theorem}

\begin{assumption}[Mode expansion]\label{ass:mode}
The state vector admits an asymptotic Oseledets decomposition $v_n = \sum_i A_i(n) u_i(n)$ with $\log|A_i(n)| = n\lambda_i + o(n\lambda_a)$, where $u_i(n)$ are the covariant Lyapunov vectors at time $n$.
\end{assumption}

\begin{assumption}[Denominator visibility]\label{ass:denvis}
$\psi(u_a(n)) \neq 0$ for all sufficiently large $n$, and $|\psi(u_a(n))|$ is not exponentially small: $\log|\psi(u_a(n))| = o(n\lambda_a)$.
\end{assumption}

\begin{assumption}[Unique dominance of $J^\star$]\label{ass:dom}
There exists $\eta > 0$ such that $\lambda_{J^\star} - \lambda_j < -\eta$ for all visible $j \neq a, J^\star$, i.e., $J^\star$ is spectrally separated from all other visible modes.
\end{assumption}

\begin{assumption}[Spectral gap]\label{ass:gap}
$\lambda_{J^\star} < \lambda_a$.
\end{assumption}

\begin{assumption}[Non-cancellation and double-depth separation]\label{ass:noncanc}
The leading visibility coefficient $V_{a,J^\star}(n)$ does not vanish identically, and the error term satisfies the double-depth separation:
\[
\frac{r_{2n} - L}{r_n - L} \to 0 \quad \text{as } n \to \infty,
\]
equivalently, $E(2n) - E(n) \to -\infty$ where $E(n) = n(\lambda_{J^\star}(n) - \lambda_a(n))$.
\end{assumption}

\begin{theorem}[Visible-mode delta]\label{thm:main}
Under Assumptions~\ref{ass:mode}--\ref{ass:noncanc}, the double-depth irrationality delta satisfies
\[
\delta_n = -\frac{\lambda_{J^\star}(n)}{\lambda_a(n)} + o(1), \qquad \delta_\infty = -\frac{\lambda_{J^\star}}{\lambda_a} \text{ when exponents converge.}
\]
\end{theorem}

\begin{proof}
\emph{Denominator.} By Assumption~\ref{ass:denvis}, $\psi(v_n) = A_a(n)\psi(u_a(n))(1 + R_n)$ where $R_n = \sum_{j \neq a} \frac{A_j(n)}{A_a(n)} \frac{\psi(u_j(n))}{\psi(u_a(n))}$. By Assumption~\ref{ass:mode}, $\log|A_j/A_a| = n(\lambda_j - \lambda_a) + o(n\lambda_a)$. For $j \neq a$ with $\lambda_j < \lambda_a$, this is $o(n\lambda_a)$. For invisible $j$ (where $\psi(u_j(n))$ is exponentially small), the contribution is also $o(n\lambda_a)$. Hence $R_n \to 0$ and $\log|q_n| = n\lambda_a + o(n\lambda_a)$.

\emph{Error.} Write $\epsilon_j(n) = A_j(n)/A_a(n)$. Then
\[
r_n - L = \frac{\sum_{j \neq a} \epsilon_j(n) V_{a,j}(\varphi, \psi; n)}{\psi(u_a(n))\left(\psi(u_a(n)) + \sum_{j \neq a} \epsilon_j(n) \psi(u_j(n))\right)}.
\]
Invisible modes drop out (their $V_{a,j}$ is exponentially small). By Assumption~\ref{ass:dom}, $J^\star$ dominates all other visible modes by a spectral gap $\eta > 0$, so $\log|r_n - L| = n(\lambda_{J^\star} - \lambda_a) + o(n\lambda_a)$.

\emph{Double depth.} By Assumption~\ref{ass:noncanc}, $E(2n) - E(n) \to -\infty$, so $r_{2n} - L$ is exponentially smaller than $r_n - L$. Therefore $r_n - r_{2n} = (r_n - L)(1 + o(1))$, and
\[
\delta_n = -1 - \frac{n(\lambda_{J^\star} - \lambda_a) + o(n\lambda_a)}{n\lambda_a + o(n\lambda_a)} = -\frac{\lambda_{J^\star}}{\lambda_a} + o(1). \qedhere
\]
\end{proof}

\begin{corollary}[Sign criterion]
If $\lambda_a > 0$, then $\delta_\infty > 0 \iff \lambda_{J^\star} < 0$.
\end{corollary}

\subsection{Exterior-product interpretation of the wedge term}
\label{sec:wedge}

The numerator $p_n q_{2n} - p_{2n} q_n$ is not, in general, a minor of the product matrix $P_n$. It is the determinant of two projective observations at times $n$ and $2n$:
\[
p_n q_{2n} - p_{2n} q_n = \det\begin{pmatrix} \varphi(v_n) & \varphi(v_{2n}) \\ \psi(v_n) & \psi(v_{2n}) \end{pmatrix} = (\varphi \wedge \psi)(v_n \wedge v_{2n}),
\]
an exterior-product observable of two transported states. In the generic case ($a=1$, $J^\star=2$), this wedge grows at scale $n(\lambda_1 + \lambda_2)$, and dividing by two denominator scales $n\lambda_1$ each leaves $n(\lambda_2 - \lambda_1)$, recovering the standard formula.

More generally, when $J^\star > 2$, the wedge is dominated by the $a$-th and $J^\star$-th modes, growing at scale $n(\lambda_a + \lambda_{J^\star})$, and the delta formula becomes $-\lambda_{J^\star}/\lambda_a$ as in Theorem~\ref{thm:main}.

\section{Failure modes and taxonomy}

\begin{center}
\begin{tabular}{lll}
\toprule
Regime & Spectral condition & Correct $\delta$ \\
\midrule
Standard positive & $a{=}1$, $J^\star{=}2$, $\lambda_2 < 0$ & $-\lambda_2/\lambda_1 > 0$ \\
Standard negative & $a{=}1$, $J^\star{=}2$, $\lambda_2 > 0$ & $< 0$ \\
Hidden positive & $V_{1,2}{=}0$, $\lambda_3 < 0$ & $-\lambda_3/\lambda_1 > 0$ \\
Denominator-blind & $a > 1$ & $-\lambda_{J^\star}/\lambda_a$ \\
Non-dominated & $\lambda_1 \approx \lambda_2$ & unstable \\
Contracting denominator & $\lambda_a < 0$ & proxy invalid \\
\bottomrule
\end{tabular}
\end{center}

The hidden-positive regime is the key failure mode of the naive $\hat\delta = -\lambda_2/\lambda_1$: a trajectory with genuinely positive $\delta$ can have $\lambda_2 > 0$ (invisible error mode) while $\lambda_3 < 0$ (the true $J^\star$). The robust detector scans the full ladder and selects the first non-degenerate ratio.

The \emph{contracting denominator} regime ($\lambda_a < 0$) is an additional failure mode discovered in the z$=\pm 1$ boundary regime: when $\lambda_1 < 0$, the ratio $-\lambda_2/\lambda_1$ can be positive even though the exact delta is negative, because the denominator is contracting rather than expanding. This motivates the $\lambda_a > 0$ filter in our screening pipeline.

\section{Classical calibration}
\label{sec:calibration}

\subsection{Ap\'ery recurrences}

\begin{center}
\begin{tabular}{lccccc}
\toprule
Constant & $n$ & $\delta_n$ (exact) & $-\lambda_2/\lambda_1$ & $|$diff$|$ \\
\midrule
$\zeta(2)$ & 100 & 1.04267 & 1.03907 & $3.6 \times 10^{-3}$ \\
$\zeta(2)$ & 1000 & 1.00610 & 1.00576 & $3.4 \times 10^{-4}$ \\
$\zeta(3)$ & 100 & $\approx 1.046$ & $\approx 1.041$ & $\sim 5 \times 10^{-3}$ \\
$\zeta(3)$ & 500 & 1.01163 & 1.01063 & $1.0 \times 10^{-3}$ \\
\bottomrule
\end{tabular}
\end{center}

Both converge to $\delta_\infty = 1 = 1/(d-1)$ for $d=2$, with the proxy tracking to $\sim 10^{-4}$ at depth 1000. The QR-based estimator $\widehat\lambda_i^{\mathrm{QR}}$ agrees with the SVD-based $\lambda_i^{\mathrm{SVD}}$ to within $10^{-3}$ relative error at $N=120$ on these systems.

\subsection{Negative-delta PCFs}

Four Ramanujan Machine PCFs ($4/(3\pi{-}8)$, $8/\pi^2$, $1/(1{-}\log 2)$, $2/(2G{-}1)$) were tested at $n=1000$. All have $\lambda_1 > \lambda_2 > 0$: the ratio converges (spectral gap) but $\delta < 0$ (error mode expands). The proxy correctly classifies all as negative.

\section{The 6F5 CMF model}

A 6F5 companion trajectory is parametrized by integer shift $\mathbf{s} \in [-\text{box},\text{box}]^{11}$, direction $\mathbf{d}$, and $z = z_{\text{num}}/z_{\text{den}}$. At step $n$, the $6 \times 6$ companion matrix $M(n)$ has sub-diagonal ones and last column built from elementary symmetric polynomials of two shifted arithmetic ladders, gauge-cleared so all entries are integers. The search space at box$=6$ is $\sim 1.5 \times 10^{15}$.

A trajectory is \emph{degenerate} if any shift component $s_i = -1$ for $i \in \{0, \ldots, d-1\}$, causing the corresponding row of the companion matrix to collapse (the $i$-th elementary symmetric polynomial vanishes). Degenerate trajectories reduce to lower-dimensional sub-systems and are filtered in production sweeps.

The prior LUMI RNS campaign screened $2.5 \times 10^9$ candidates in $\sim$20{,}000 GPU-hours, yielding $\sim$300 confirmed positives. The estimated positive base rate was $\sim 1.2 \times 10^{-7}$.

\section{GPU spectral algorithm}

The proxy $\hat\delta = -\lambda_2(N)/\lambda_1(N)$ is computed via QR decomposition at depth $N=120$. Cost: $O(d^3 N)$ flops with fixed-width floating-point numbers. On LUMI (AMD MI250X), a custom QR kernel achieves $\sim$301{,}000 trajectories/s per GCD (Graphics Compute Die). Each MI250X module contains two GCDs, each with 64\,GB HBM2e memory; LUMI bills one GPU-hour per MI250X module (two GCDs) \cite{lumi_hardware,lumi_billing}. Throughput is therefore $\sim 1.08 \times 10^9$ trajectories per GCD-hour, or $\sim 2.17 \times 10^9$ per billed GPU-hour.

The full-spectrum ladder $r_k = -\lambda_k/\lambda_1$ for $k = 2, \ldots, d$ is extracted simultaneously at no additional cost.

For blind screening, the conservative full-ladder detector keeps a candidate if \emph{any} stable $r_k > 0$ (with $|r_k| < 1.05$). This catches hidden-positive regimes where $r_2 < 0$ but $r_3 > 0$. Two additional filters are applied:
\begin{itemize}
\item \emph{Positive denominator filter} ($\lambda_1 > 0$): rejects trajectories with contracting denominators, which produce false-positive proxy values.
\item \emph{Non-degenerate filter}: rejects trajectories where any $s_i = -1$ for $i \in \{0, \ldots, d-1\}$.
\end{itemize}

\begin{algorithm}
\caption{GPU proxy screening (Stage 1)}
\begin{algorithmic}[1]
\State \textbf{Input:} seed, gid range $[0, N)$, depth $N_{\text{steps}}$, box, dim $d$
\State Build direction pool and $z$-pool on device
\For{each batch of $B$ trajectories}
  \State Generate parameters $(\text{shift}, \text{dir}, z)$ from (seed, gid) on device
  \State For each trajectory: compute $M(n)$ for $n = 1, \ldots, N_{\text{steps}}$
  \State Per-step QR: accumulate $\log|R_{k,ii}|$ for $i = 1, \ldots, d$
  \State Compute $r_k = -\widehat\lambda_k / \widehat\lambda_1$ for $k = 2, \ldots, d$
  \State Apply filters: $\widehat\lambda_1 > 0$, non-degenerate, $r_k > \text{thresh}$
  \State Write survivors (gid, $r_k$, $\widehat\lambda_1$, shift, dir, $z$) to JSONL
\EndFor
\end{algorithmic}
\end{algorithm}

\section{Benchmark against exact arithmetic}
\label{sec:benchmark}

\subsection{Precision comparison}

The proxy supports both float32 and float64 QR. The production LUMI sweeps use float32 for throughput; validation against exact arithmetic uses float64. Table~\ref{tab:precision} reports the precision-dependent benchmark.

\begin{table}[h]
\centering
\caption{Precision comparison: proxy vs exact on 2000 matched candidates (dim=6, $N=120$, box=6, CPU 8-core)}
\label{tab:precision}
\begin{tabular}{lcccc}
\toprule
Precision & Throughput (traj/s) & Sign errors & MAE & Notes \\
\midrule
float64 & 1{,}192 & 0 / 2000 & 0.013 & All negatives; no positives in sample \\
float32 & $\sim$2{,}400 & 0 / 2000 & $\sim$0.02 & GPU production uses float32 \\
Exact (RNS) & 909 & --- & --- & Integer arithmetic reference \\
\bottomrule
\end{tabular}
\end{table}

\begin{remark}
The CPU benchmark sample contains no positive-delta trajectories (consistent with base rate $\sim 10^{-6}$ at box=6). Sign agreement is therefore computed on negatives only. Precision, recall, and cost-per-hit are \emph{not estimable} from this sample ($\text{TP} = 0$). Stratified validation with positive cases is reported in Section~\ref{sec:validation}.
\end{remark}

\subsection{GPU throughput}

\begin{center}
\begin{tabular}{lcc}
\toprule
Metric & Proxy (float32 QR) & Exact (RNS integer) \\
\midrule
Throughput per GCD-hour & $\sim 1.08 \times 10^9$ & $\sim 10^5$ \\
Raw speedup & \multicolumn{2}{c}{$\sim 10^4 \times$} \\
Memory per trajectory & $O(d^2)$ fixed & $O(d^2 \cdot N \log N)$ growing \\
Scaling with depth $N$ & $O(N)$ & $O(N^2 \log N)$ bit-cost \\
\bottomrule
\end{tabular}
\end{center}

\textbf{Benchmark protocol:} The reproducible benchmark script (\texttt{benchmark\_proxy\_vs\_exact.py} \cite{deltaproxy_repo}) measures four distinct metrics on identical candidates at matched depth:

\begin{enumerate}
\item \emph{Raw throughput:} trajectories/second for proxy vs exact.
\item \emph{Statistical enrichment:} precision, recall, sign-agreement of proxy screening.
\item \emph{End-to-end discovery speed-up:} GPU-time per exact-verified positive.
\item \emph{Verification cost:} exact arithmetic time on the selected tail.
\end{enumerate}

When no positives appear in the benchmark sample ($\text{TP} = 0$), the script reports \texttt{null} for cost-per-hit and precision/recall, rather than dividing by zero. The sign agreement and MAE remain valid as negative-only metrics.

\section{Accuracy and reliability}
\label{sec:validation}

\subsection{Controlled calibration}

On a 1-D calibration band ($z = k/20$, $N=180$): sign accuracy 1.000, MAE 0.0088, Pearson $r = 0.9999$.

\subsection{Independent holdout validation}

Five structurally disjoint holdout sets (500 trajectories each, $N{=}120$, box${=}6$) were evaluated (\texttt{validation\_holdout.py} \cite{deltaproxy_repo}):

\begin{center}
\begin{tabular}{lcccccc}
\toprule
Holdout & $n$ & sign & MAE & $r$ & prec & rec \\
\midrule
Shard (seed=99999) & 500 & 1.0000 & 0.0125 & 0.989 & --- & --- \\
Direction & 500 & 1.0000 & 0.0097 & 0.999 & --- & --- \\
$z$-holdout & 500 & 1.0000 & 0.0130 & 0.986 & --- & --- \\
Full-dim only & 500 & 1.0000 & 0.0137 & 0.983 & --- & --- \\
Adversarial & 500 & 1.0000 & 0.0289 & 0.962 & --- & --- \\
\bottomrule
\end{tabular}
\end{center}

All five holdouts achieve perfect sign agreement. The adversarial set (spectral gap $\in [0.25, 2)$) shows higher MAE (0.029 vs $\sim 0.01$) but no sign errors. No positives appeared in these random box-6 samples (consistent with base rate $\sim 10^{-6}$), so precision/recall are not reported here; they are measured in the billion-scale sweeps (Section~\ref{sec:sweeps}).

\subsection{Stratified validation with positive cases}

To address the limitation that random samples contain no positives, we construct a stratified test set including known positive-delta trajectories:

\begin{center}
\begin{tabular}{lcccc}
\toprule
Stratum & $n$ & Sign agree & MAE & Notes \\
\midrule
Generic negatives & 2000 & 1.000 & 0.013 & Random box-6 \\
Ap\'ery $\zeta(2)$ & 1 & 1.000 & 0.004 & $\delta = 1$, $d=2$ \\
Ap\'ery $\zeta(3)$ & 1 & 1.000 & 0.005 & $\delta = 1$, $d=2$ \\
Degenerate 3F2 (Hit A) & 1 & 1.000 & 0.006 & $\delta \to 0.5$, $d=3$ effective \\
Non-degenerate 6F5 (Hit B) & 1 & 1.000 & 0.015 & $\delta \to 0.15$, $d=6$ \\
Contracting denominator & 10 & 0.80 & 0.31 & $\lambda_1 < 0$; proxy fails \\
\bottomrule
\end{tabular}
\end{center}

The contracting-denominator stratum confirms the failure mode identified in Section~5: when $\lambda_1 < 0$, the proxy $-\lambda_2/\lambda_1$ can be positive while exact $\delta$ is negative. The $\lambda_1 > 0$ filter eliminates these false positives.

\subsection{Claim: proxy is a conjecture generator, not a certificate}

The proxy is a very effective screening statistic in dominated regimes. It is \emph{not} a universal arithmetic certificate. Every arithmetic claim in this paper is settled by exact integer computation.

\section{Billion-scale screening experiment}
\label{sec:sweeps}

\subsection{LUMI 6F5 sweeps}

Two independent 200-billion-trajectory GPU sweeps (seeds 20260626 and 20270704) produced 559{,}605 and 559{,}530 proxy survivors ($\hat\delta > 0.02$). Exact verification of 559{,}417 new survivors confirmed 551{,}149 with $\delta > 0$ (98.5\%). Proxy vs exact: sign-agree = 0.996, MAE = 0.0497.

The two runs agree to $0.10\sigma$ of the binomial mean, confirming a reproducible empirical hit rate $p \approx 2.8 \times 10^{-6}$ under the specified generator and sampling measure.

\subsection{Local billion-trajectory sweep}

A local sweep of $10^9$ trajectories at $z = \pm 1/2$ (6 CPU workers, 13.5 hours) produced 104{,}543 proxy survivors, of which 52{,}546 had exact $\delta > 0$ (50.2\% confirmation). PSLQ identified 70 hits with integer relations; 12 remained unidentified by PSLQ and are classified as candidate non-rational limits pending further analysis.

\subsection{z$=\pm 1$ boundary sweep}

A dedicated sweep restricted to $z = +1$ and $z = -1$ (the boundary regime where the gauge factor equals $\pm 1$) was conducted to probe the positive-delta tail. Local testing (200{,}000 trajectories, seed 20260717) with the $\lambda_1 > 0$ and non-degenerate filters active yielded:

\begin{itemize}
\item Positive proxy rate: $\sim 5 \times 10^{-6}$ (1 survivor in 200k)
\item All strong proxy values ($r_2 > 0.5$) are false positives from degenerate or contracting-denominator trajectories
\item The single non-degenerate survivor has $r_2 = +0.001$ with exact $\delta = +0.004$ at $N=400$ (weak positive, $\delta \to 0$)
\end{itemize}

A production LUMI sweep of 16 waves $\times$ 8 nodes $\times$ 24h (24{,}576 GPU-hours) is currently screening $\sim 1.33 \times 10^{13}$ trajectories (39\% of the $6.81 \times 10^{13}$ total $z = \pm 1$ space).

\subsection{Multi-ratio spectral diagnostic}

Full Lyapunov ladders across 4F3/5F4/6F5/8F7 (40{,}000 samples each, depth 200):

\begin{center}
\begin{tabular}{llcccc}
\toprule
Family & Matrix & $r_2 > 0.02$ & $r_2$ max & $r_3 > 0.02$ & Deeper signal \\
\midrule
4F3 & pFq & 46 & $+2.698$ & 7{,}465 & 16{,}311 \\
5F4 & pFq & 0 & $-0.068$ & 13 & 21{,}640 \\
6F5 & pFq & 0 & $-0.136$ & 0 & 17{,}356 \\
8F7 & pFq & 0 & $-0.591$ & 0 & 0 \\
6F5 & companion & 0 & $+0.008$ & 27 & 21{,}971 \\
\bottomrule
\end{tabular}
\end{center}

The spectral signal $r_2 > 0.02$ collapses monotonically with dimension: as $d$ grows, the cocycle typically develops multiple positive Lyapunov exponents. The deeper-ratio columns ($r_3 > 0.02$ and total deeper signal) should be interpreted as \emph{spectral candidates}, not certified signals: positive $r_k$ for $k \geq 3$ requires visibility analysis and exact arithmetic verification to confirm.

\section{Experimental geometry}

\subsection{Confirmed hits}

\paragraph{Hit A (degenerate 3F2 in 6F5, $\delta \to +0.50$).}
Shift $= [-1,-1,-1,0,0,0,-2,-2,-2,-2,3]$, $z=1/3$. Rows 0--2 collapse, yielding a 3D CMF with $a_{n+3} = 3(n{+}4)a_{n+2} - 3a_{n+1} - a_n$. Delta: $n=1000$: $+0.494$, converging to $1/(3{-}1) = 0.5$. Limit $L \approx 0.32722\ldots$ (unidentified by PSLQ). This is an asymptotic positive-delta candidate; a complete irrationality proof is not included in this paper.

\paragraph{Hit B (non-degenerate 6F5, $\delta \to +0.15$, increasing).}
{\sloppy Shift $= [-2,-2,0,0,-2,0,-2,-2,-2,0,-2]$, $z=7/20$. All six rows active ($f_i \neq 0$). Recurrence: $a_{n+6} = 20(n{+}2)a_{n+5} + 21a_{n+4} - 21a_{n+2} + 7a_n$.\par}

Exact $\delta = 0.1432$ at $N{=}120$. The finite-depth arithmetic delta is numerically closest to the fourth Lyapunov ratio: $r_4 = -\lambda_4/\lambda_1 = 0.146$ matches exact $\delta$ to within 0.003, while the naive $r_2 = 0.128$ underpredicts by 0.015. This is consistent with the hidden-visible-mode regime, in which $\lambda_2$ and $\lambda_3$ are not the dominant error modes for this readout. We note that this observation is based on post-hoc comparison with the known exact delta; independent identification of $J^\star = 4$ via visibility analysis (covariant Lyapunov vectors or explicit SVD of $P_N$) is left for future work. Delta: $n=400$: $+0.149$ (toward $1/5$).

\subsection{Delta Ladder}

\begin{conjecture}[Delta Ladder]
For a $d$-dimensional CMF with full effective dimension,
\[
\delta_\infty \approx \frac{1}{d_{\mathrm{eff}} - 1}.
\]
\end{conjecture}

This is supported by: $d=2$ (Ap\'ery): $\delta = 1 = 1/(2-1)$. $d=3$ (Hit A): $\delta \to 0.5 = 1/(3-1)$. $d=6$ (Hit B): $\delta \to \sim 0.15 \approx 1/(6-1) = 0.2$ (increasing). This remains a conjecture, not a theorem.

\section{Practical \texorpdfstring{$J^\star$}{J*} estimation}
\label{sec:jstar}

We define three modes of $J^\star$ estimation, with distinct epistemic status:

\begin{enumerate}
\item \textbf{Blind full-ladder screening} (\texttt{jstar\_estimator.py --mode blind}): keeps a trajectory if \emph{any} $-\lambda_k/\lambda_1 > 0$. This is a broad candidate-generation filter. The selected $k_{\mathrm{candidate}} = \arg\max_k (-\lambda_k/\lambda_1)$ should not be called $J^\star$; it is $k_{\mathrm{max\text{-}positive}}$, a high-recall but low-precision screening statistic. In high dimensions, deep Lyapunov modes are often negative without affecting the observable.

\item \textbf{Post-hoc analysis} (\texttt{jstar\_estimator.py --mode posthoc}): computes exact $\delta$ and selects the Lyapunov ratio closest to it. This is useful for exploratory analysis but is \emph{circular} as an estimator: it uses the target value to select the model. For Hit B, the observation that $r_4$ is closest to exact $\delta$ should be read as: ``the finite-depth arithmetic delta is numerically closest to the fourth Lyapunov ratio,'' not as an independent identification of $J^\star = 4$.

\item \textbf{Singular-vector method} (\texttt{jstar\_estimator.py --mode sv}): attempts to estimate visibility from the columns of the accumulated $Q$ matrix. The current implementation uses $\varphi = e_1^\top$, $\psi = e_d^\top$, which may not match the actual readout pair producing $p_n/q_n$. A correct implementation requires either explicit SVD of $P_N$, backward-forward covariant Lyapunov vector computation, or adjoint mode propagation with the actual functionals $\varphi, \psi$. This is left for future work.
\end{enumerate}

\section{Reproducibility and certificates}

\subsection{Verification architecture}

The two-stage funnel separates screening from certification:

\begin{enumerate}
\item \textbf{Stage 1 (GPU proxy):} Screen billions; keep $\hat\delta > \text{thresh}$ survivors. Parameters re-derivable from (seed, gid), so only survivors are written.
\item \textbf{Stage 2 (exact verify):} Pure-integer double-depth walk confirms each survivor. Independent of proxy code.
\end{enumerate}

The verification ladder: V0 (float flag) $\to$ V1 (independent integer engine) $\to$ V2 (exact confirmation) $\to$ V3 (CAS certificates: mpmath, pure-integer Python, SageMath) $\to$ V4 (monotonic convergence) $\to$ V5 (formal theorem, open).

\subsection{Reproducibility package}

All code is available at \url{https://github.com/VesterlundCoder/deltaproxy} \cite{deltaproxy_repo} under the MIT license. The repository includes:

\begin{itemize}
\item \texttt{benchmark\_proxy\_vs\_exact.py}: reproducible benchmark with environment logging, handles TP=0 correctly.
\item \texttt{validation\_holdout.py}: five independent holdout validations.
\item \texttt{jstar\_estimator.py}: blind, post-hoc, and singular-vector $J^\star$ estimation.
\item \texttt{gpu\_proxy/}: GPU kernel (\texttt{proxy\_kernel.py}), sweep driver (\texttt{gpu\_sweep.py}), parameter generation (\texttt{params.py}), calibration (\texttt{calibrate.py}).
\item \texttt{cmf\_generic.py}: CMF construction and exact arithmetic delta.
\item \texttt{spectral\_delta.py}: Lyapunov spectrum computation.
\item \texttt{pslq\_companion.py}: PSLQ integer relation identification.
\item \texttt{verify\_survivors.py}: Stage-2 exact verification.
\item \texttt{deploy\_zpm1\_full.sh}: LUMI deployment script (rsync + SLURM submission).
\item SLURM batch scripts: \texttt{lumi\_zpm1\_positive\_669.sbatch}, \texttt{lumi\_sweep\_669.sbatch}.
\item Random seeds: 20260626 (sweep 1), 20270704 (sweep 2), 99999 (shard holdout), 20260717 (z$=\pm 1$ sweep).
\item LUMI container: PyTorch/2.6.0-rocm-6.2.4 (SIF path in SLURM scripts).
\end{itemize}

\subsection{Data availability}

Survivor manifests and verification certificates from the 200-billion-trajectory sweeps are available in the repository. The z$=\pm 1$ production sweep results will be deposited upon job completion. All trajectories are re-derivable from (seed, gid) using the deterministic parameter generator in \texttt{params.py}, which is bit-identical to the on-device GPU generator.

\section{Claim matrix}

\begin{center}
\begin{tabular}{ll}
\toprule
Class & Claim \\
\midrule
\textbf{Theorem} & Visible-mode delta theorem (Thm.~\ref{thm:main}) \\
\textbf{Exact identity} & Exterior-product formula for $r_n - r_{2n}$ \\
\textbf{Verified computation} & Exact arithmetic for specified trajectories \\
\textbf{Empirical law} & Proxy tracks exact delta in test population \\
\textbf{Benchmark result} & $\sim 10^9$ trajectories/GCD-hour, $\sim 10^4$ speedup \\
\textbf{Empirical rate} & Reproducible hit rate $p \approx 2.8 \times 10^{-6}$ \\
\textbf{Conjecture} & Generic $J^\star$-CMF law \\
\textbf{Conjecture} & Delta Ladder: $\delta \to 1/(d_{\mathrm{eff}}-1)$ \\
\textbf{Spectral candidate} & Deeper-ratio signals ($r_3, r_4, \ldots$) \\
\textbf{Research program} & Global CMF/Stokes geometry \\
\bottomrule
\end{tabular}
\end{center}

\section{Discussion and limitations}

\subsection{What is proven}
The visible-mode delta theorem (Thm.~\ref{thm:main}) is a rigorous conditional result under five explicit assumptions. The exterior-product identity is exact. The Ap\'ery calibration is verified to $\sim 10^{-4}$.

\subsection{What is numerical}
The proxy accuracy (sign agreement $\sim 99\%$, MAE $\sim 0.03$) is an empirical observation on test populations. The $\sim 10^4$ speedup is a measured benchmark. The hit rate $p \approx 2.8 \times 10^{-6}$ is a reproducible empirical rate under the specified sampling measure, not a mathematical invariant.

\subsection{What is conjectural}
The Delta Ladder is a conjecture supported by three data points. The generic $J^\star$-CMF law (that $J^\star = 2$ generically) is a conjecture. Deeper-ratio signals ($r_3, r_4, \ldots$) are spectral candidates, not certified signals, until visibility analysis and exact arithmetic confirm them.

\subsection{What is not claimed}
\begin{itemize}
\item Hit A is an \emph{asymptotic positive-delta candidate}, not a proven irrational number. A complete irrationality proof is not included.
\item PSLQ non-identification does not certify non-rationality; it reports the absence of detected integer relations at the tested precision.
\item The post-hoc $J^\star = 4$ observation for Hit B is not an independent identification; it is a numerical coincidence between the known exact delta and a Lyapunov ratio.
\item The proxy is not a universal arithmetic certificate; it fails in contracting-denominator and non-dominated regimes.
\end{itemize}

\subsection{Future directions}
\begin{itemize}
\item Independent $J^\star$ identification via covariant Lyapunov vectors or explicit SVD of $P_N$.
\item 7F6/8F7 production sweeps with the non-degenerate and $\lambda_a > 0$ filters.
\item Model-guided sampling and full spectral atlases.
\item Complete irrationality proofs for Hit A and Hit B.
\item Broader literature comparison and theoretical connections to Oseledets multiplicative ergodic theory.
\end{itemize}

\section{Conclusion}

One law organizes the subject: $\delta = -\lambda_{J^\star}/\lambda_a$. Projective convergence requires a spectral gap; positive delta requires a visible contracting mode. The law is exact enough to calibrate on Ap\'ery ($10^{-4}$), sharp enough to classify PCFs, and cheap enough to accelerate 6F5 CMF discovery by $10^4$---provided its hypotheses are diagnosed via the full Lyapunov ladder, gauge conventions are specified, and every arithmetic claim is settled by exact integer computation.

\section*{Data and software availability}

All code, benchmark scripts, and reproducibility documentation are available at \url{https://github.com/VesterlundCoder/deltaproxy} (MIT license, version-controlled). The repository contains the GPU proxy kernel, sweep drivers, exact verification code, PSLQ identification, benchmark and validation scripts, and LUMI SLURM deployment scripts. All trajectories are re-derivable from deterministic seeds.

\section*{Author contributions}

D.V.\ designed the study, developed the theory, implemented all code, conducted all experiments, and wrote the manuscript.

\begin{thebibliography}{9}
\bibitem{apery1979} R.~Ap\'ery, Irrationalit\'e de $\zeta(2)$ et $\zeta(3)$, Ast\'erisque \textbf{61} (1979), 11--13.
\bibitem{vanderpoorten1979} A.~van der Poorten, A proof that Euler missed..., Math. Intelligencer \textbf{1} (1979), 195--203.
\bibitem{raayoni2021} G.~Raayoni et al., Generating conjectures on fundamental constants with the Ramanujan Machine, Nature \textbf{590} (2021), 67--73.
\bibitem{elimelech2024} R.~Elimelech et al., Algorithm-assisted discovery of an intrinsic order among mathematical constants, PNAS \textbf{121} (2024), e2321440121.
\bibitem{oseledets1968} V.~I.~Oseledets, A multiplicative ergodic theorem, Trans. Moscow Math. Soc. \textbf{19} (1968), 197--231.
\bibitem{ferguson1999} H.~R.~P.~Ferguson, D.~H.~Bailey, S.~Arno, Analysis of PSLQ, Math. Comp. \textbf{68} (1999), 351--369.
\bibitem{rabinowitz1991} S.~Rabinowitz, A survey of finite continued fractions, in \emph{Computational aspects of number theory}, 1997.
\bibitem{borwein1992} J.~M.~Borwein, P.~B.~Borwein, \emph{Pi and the AGM}, Wiley, 1987.
\bibitem{zudilin2001} W.~Zudilin, One of the numbers $\zeta(5), \zeta(7), \zeta(9), \zeta(11)$ is irrational, Russ.~Math.~Surv.~\textbf{56} (2001), 774--776.
\bibitem{krattenthaler1997} C.~Krattenthaler, H.~M.~Srivastava, Summations for basic hypergeometric series, J.~Math.~Anal.~Appl.~\textbf{207} (1997), 12--27.
\bibitem{lumi_hardware} LUMI Supercomputer, GPU nodes documentation, \url{https://docs.lumi-supercomputer.eu/hardware/lumig/}, 2024.
\bibitem{lumi_billing} LUMI Supercomputer, Billing policy, \url{https://docs.lumi-supercomputer.eu/runjobs/lumi_env/billing/}, 2024.
\bibitem{deltaproxy_repo} D.~Vesterlund, Delta Proxy: Visible-Mode Lyapunov Geometry and Spectral Screening for CMFs, \url{https://github.com/VesterlundCoder/deltaproxy}, 2026.
\bibitem{vesterlund2026} D.~Vesterlund, Spectral dynamics report, project\_465002669 (LUMI), 2026. Available at \cite{deltaproxy_repo}.
\end{thebibliography}

\end{document}
